\documentclass[paper=screen,mode=present,style=zysimple]{powerdot}

\usepackage{fancyvrb,verbatim,listings,bbold}
\usepackage{proof}
\usepackage[all]{xypic,xy}

\usepackage{latexsym,stmaryrd,amsmath,amsfonts}

%% \newtheorem{definition}{Definition}
%% \newtheorem{theorem}{Theorem}
%% \newtheorem{lemma}{Lemma}
\newtheorem{prop}{Proposition}
%% \newtheorem{corollary}{Corollary}
%% \newenvironment{proof}{{\sc Proof . }}{{\hfill $\square$}}
%% \newcommand{\myqed}{\hfill $\square$}
%% \newcommand{\ped}{\vspace{-1.6em}}
%% \newcommand{\note}[1]{\noindent\fbox{\parbox{\textwidth}{\em #1}}}

\newcommand{\abind}[2]{\mathtt{bind} \ (\, {#1} \,)  \ \mathtt{in} \ {#2}}
\newcommand{\acase}{\mathtt{case}}
\newcommand{\adv}{\mathcal{A}}
\newcommand{\aspecto}[1]{\xrightarrow{\; {#1} \;}}
\newcommand{\asplinear}{\mathfrak{l}}
\newcommand{\aspmodal}{\mathfrak{m}}
\newcommand{\aspnormal}{\mathfrak{n}}
%
\newcommand{\betareduce}{\leadsto_\beta}
\newcommand{\betareduceto}{\twoheadrightarrow_\beta}
\newcommand{\bit}{\mathtt{B}}
\newcommand{\bitone}{\mathtt{B}_1}
\newcommand{\bitzero}{\mathtt{B}_0}
%
\newcommand{\caseto}{\triangleright}
\newcommand{\cbind}[3]{\mathtt{bind} \ {#1} = {#2} \ \mathtt{in} \ {#3}}
\newcommand{\ccase}[4]{\mathtt{case} \; {#1} \; \mathtt{of} \; \cnil \caseto {#2}; \;
                                          \bitzero x_0 \caseto {#3}; \; \bitone x_1 \caseto {#4}}
\newcommand{\cequiv}{\simeq}
\newcommand{\cif}[3]{\mathtt{if} \; {#1} \; \mathtt{then} \; {#2} \; \mathtt{else} \; {#3}}
\newcommand{\ctrue}{\mathtt{true}}
\newcommand{\cfalse}{\mathtt{false}}
\newcommand{\cfix}{\mathtt{fix}\,}
\newcommand{\cfixin}[2]{\mathtt{fix}\; {#1}\; \mathtt{in} \; {#2}}
\newcommand{\cfold}{\mathtt{fold}\,}
\newcommand{\chd}{\mathtt{hd}}
\newcommand{\cinj}{\mathtt{inj}}
\newcommand{\clift}[1]{\lfloor{#1}\rfloor}
\newcommand{\clet}[3]{\mathtt{let} \ {#1} = {#2} \ \mathtt{in} \ {#3}}
\newcommand{\cletst}[3]{\mathtt{liftlet} \ {#1} = {#2} \ \mathtt{in} \ {#3}}
\newcommand{\cletrec}[3]{\mathtt{letrec} \ {#1} = {#2} \ \mathtt{in} \ {#3}}
\newcommand{\cnil}{\mathtt{nil}}
\newcommand{\conc}{\! \bullet \!}
\newcommand{\concat}{\pmb{\mathit{conc}}}
\newcommand{\context}{\mathbb{C}}
\newcommand{\cop}{\;\mathtt{op}\;}
\newcommand{\cpair}[1]{\langle{#1}\rangle}
\newcommand{\cproj}{\mathtt{proj}}
\newcommand{\crand}{\mathtt{rand}}
\newcommand{\crec}{\mathtt{rec}}
\newcommand{\croll}{\mathtt{roll}\,}
\newcommand{\crule}[1]{\mbox{{\textsl{#1}}}}
\newcommand{\cslrtype}{\Vdash}
\newcommand{\csucc}{\mathtt{S}}
\newcommand{\ctl}{\mathtt{tl}}
\newcommand{\cut}{\pmb{\mathit{cut}}}
\newcommand{\cval}{\mathtt{return}}
\newcommand{\cwhile}[2]{\mathtt{while}\; {#1} \; \mathtt{do} \; {#2} \; \mathtt{od}}
\newcommand{\cvalof}{\pmb{valueof}}
\newcommand{\cunfold}{\mathtt{unfold}\,}
\newcommand{\cunit}{\mathtt{unit}}
\newcommand{\cunroll}{\mathtt{unroll}\,}
\newcommand{\cupdate}{\pmb{update}}
\newcommand{\czero}{\mathtt{O}}
%
\newcommand{\deduces}{\vdash}
\newcommand{\defeq}{\stackrel{\mathrm{def}}{=}}
\newcommand{\delay}[1]{\left\lfloor{#1}\right\rfloor}
\newcommand{\distr}{\mathcal{D}}
%
\newcommand{\eagerreduce}{\xrightarrow{\mathrm{eager}}}
\newcommand{\Eagerreduce}{\stackrel{\mathrm{eager}}{\Reduce}}
\newcommand{\eqtest}{\stackrel{\mbox{\tiny ?}}{=}}
\newcommand{\Eval}[1]{{\left\llbracket{#1}\right\rrbracket}}
\newcommand{\evalcontext}{\mathbb{E}}
\newcommand{\Example}{\emgreen{\pmb{\textsc{Example}\;}}}
\newcommand{\expansion}[1]{\tilde{{#1}}}
%
\newcommand{\ford}{\mathrm{ord}}
%
\newcommand{\getval}{\leftarrow}
%
\newcommand{\head}{\pmb{\mathit{hd}}}
\newcommand{\termhole}{[\;]}
\newcommand{\hoaretriple}[3]{\{{#1}\} \; {#2} \; \{{#3}\}}
%
\newcommand{\imply}{\Rightarrow}
\newcommand{\indist}{\simeq}
\newcommand{\iteration}{\pmb{\mathit{it}}}
%
\newcommand{\lazyreduce}{\xrightarrow{\mathrm{lazy}}}
\newcommand{\Lazyreduce}{\lazyreduce\!\!\!\!\!\to}
\newcommand{\ldot}{\, . \, }
\newcommand{\linearto}{\multimap}
\newcommand{\logrel}{(\mathcal{R}_\tau)_{\mbox{$\tau$ type}}}
\newcommand{\lr}[1]{\mathcal{R}_{#1}}
\newcommand{\lub}{\sqcup}
%
\newcommand{\metalambda}{\underline{\lambda}}
\newcommand{\modal}{\square}
%
\newcommand{\nat}{\mathbb{N}}
\newcommand{\nReduce}{\nrightarrow\!\!\!\!\to}
\newcommand{\nreduceto}{\nrightsquigarrow}
\newcommand{\nullenv}{\varnothing}
%
\newcommand{\orc}{\mathcal{O}}
\newcommand{\oslrtype}{\Vdash}
%
\newcommand{\porder}{\sqsubseteq}
\newcommand{\prefix}{\pmb{\mathit{pref}}}
\newcommand{\proba}{\mathbf{Pr}}
%
\newcommand{\randlen}{\pmb{\mathit{rl}}}
\newcommand{\randnum}{\pmb{\mathit{rn}}}
\newcommand{\randstr}{\pmb{\mathit{rs}}}
\newcommand{\reduce}{\leadsto}
\newcommand{\reduceto}{\rightsquigarrow\hspace*{-6px}\shortrightarrow}
\newcommand{\Reduce}{\to\!\!\!\to}
\newcommand{\Reduceto}{\Downarrow}
%
\newcommand{\set}[1]{\{{#1}\}}
\newcommand{\setbits}{\mathbb{B}}
\newcommand{\slrtype}{\vdash}
\newcommand{\splt}{\pmb{\mathit{split}}}
\newcommand{\suffix}{\pmb{\mathit{suff}}}
\newcommand{\sysname}{\pmb{\mathcal{C}}}
%
\newcommand{\tail}{\pmb{\mathit{tl}}}
\newcommand{\tbits}{\mathsf{Bits}}
\newcommand{\tbool}{\mathsf{Bool}}
\newcommand{\tnat}{\mathsf{Nat}}
\newcommand{\tcomp}{\mathsf{T}}
\newcommand{\tlift}[1]{\lfloor{#1}\rfloor}
\newcommand{\tlist}{\mathsf{List}}
\newcommand{\tself}{\mathsf{self}}
\newcommand{\tstore}{\mathsf{Store}}
\newcommand{\tstream}{\mathsf{Stream}}
\newcommand{\tloc}{\mathsf{Loc}}
\newcommand{\ttcase}{\mathtt{case}}
\newcommand{\ttif}{\mathtt{if}}
\newcommand{\ttthen}{\mathtt{then}}
\newcommand{\ttelse}{\mathtt{else}}
\newcommand{\ttin}{\mathtt{in}}
\newcommand{\ttlet}{\mathtt{let}}
\newcommand{\ttof}{\mathtt{of}}
\newcommand{\tunit}{\mathsf{Unit}}
\newcommand{\Type}{\mathcal{T}\!\mathit{ype}}
%
\newcommand{\vbool}{\pmb{\mathbb{B}}}
\newcommand{\vff}{\mathfrak{f}\!\mathfrak{f}}
\newcommand{\vnat}{\pmb{\mathrm{N}}}
\newcommand{\vnil}{\epsilon}
\newcommand{\vtt}{\mathfrak{t}\!\mathfrak{t}}
%
\newcommand{\weakpre}{\pmb{wp}}




\newcommand{\cate}{\mathcal C}

\pdsetup{lf={Theory and Practice of Functional Programming}}

\title{{\small \textsc{Theory and Practice of Functional Programming}}
\\[12pt]
{\Large \color{TitleColor} Denotations, Categories and Monads}
}
\author{\small \textbf{Dr. ZHANG Yu}
  \\[2pt] \small Institute of Software, Chinese Academy of Sciences 
  \\[5pt] \includegraphics[height=.04\slideheight]{iscas.eps} 
  \qquad \qquad 
  \includegraphics[height=.04\slideheight]{cas.eps}
}
\date{\scriptsize Fall term, 2011 \qquad GUCAS, Beijing}

\newcommand{\RED}{\mathsf{RED}}

\begin{document}

\maketitle 

\begin{slide}{Denotations of programs}
\begin{itemize}
\item All programs in same line denote the same objects:
\begin{itemize}\small
\item $6$, $2 + 4$, $3 * 2$, ...
\item $\ctrue$, $2 \eqtest 1 + 1$, $\cproj_1 \cpair{\ctrue, 0}$, ...
\item $\lambda x \ldot x * 2$, $\lambda y \ldot y + y$, ...
\item Various Haskell implementations of fibonacci function that we have seen.
\end{itemize}
\item A \emblue{denotational model (semantics)} assigns a (mathematical) object to each program. 
\begin{itemize}
\item Programs having the same denotation are considered as ``equivalent''.
\item A naive model for typed $\lambda$-calculi is based on \emred{sets}.
\end{itemize}
\end{itemize}
\end{slide}


\begin{slide}{The set model for typed $\lambda$-calculi}
\begin{itemize}
\item The \emblue{set model} is a very basic denotational model
% \vspace*{-0.5em}
\item \emblue{Types} are denoted by \emblue{sets}
% \vspace*{-0.6em}
\begin{smalltext}
\[
\begin{array}{l}
\Eval \tbool = \set{\vtt, \vff}, \quad 
\Eval \tnat = \set{0, 1, 2, \ldots } 
\\[0.3em]
\Eval{\tau_1 \times \tau_2} = \set{(v_1, v_2) \mid v_1 \in \Eval{\tau_1}, v_2 \in \Eval{\tau_2}} 
\\[0.3em]
\Eval{\tau_1 + \tau_2} = \set{(\mathbb{1}, v) \mid v \in \Eval{\tau_1}} \cup \set{(\mathbb 2, v) \mid v \in \Eval{\tau_2}} 
\\[0.3em]
\Eval{\tau_1 \to \tau_2} = \set{f \mid \empurple{\mbox{$f$ is a function from $\Eval{\tau_1}$ to $\Eval{\tau_2}$}}}
\end{array}
\]
\end{smalltext}
\vspace*{-1em}
\item \emblue{Programs} are interpreted as \emblue{functions} between sets
\begin{itemize} \small
\item $\Eval{x_1: \tau_1, \ldots, x_n: \tau_n \vdash e : \tau}$ is 
  \empurple{a function from $\Eval{\tau_1} \times \cdots \times \Eval{\tau_n}$ to $\Eval \tau$}. 
  \\
  A closed term is a constant function that returns a fixed value of $\Eval \tau$.
\item An \emblue{environment} (w.r.t. a typing context $\Gamma = \set{x_1 : \tau_1, \ldots, x_n : \tau_n}$) 
  is a function that maps $x_1$ to an element in $\tau_i$.
\item Given an environment $\rho$, \emblue{$\Eval e_\rho$} is the value comptued by applying the function  
  $\Eval {\Gamma \vdash e: \tau}$ to the tuple $(\rho(x_1), \ldots, \rho(x_n))$, and $\Eval e_\rho \in \Eval \tau$.
% \vspace*{-0.6em}
\end{itemize}
\end{itemize}
\end{slide}


\begin{slide}{The set model for typed $\lambda$-calculi}
\begin{itemize}
\item Denotations of programs
\[\small
\begin{array}{l}
\Eval{\ctrue}_\rho = \vtt, 
\Eval{0}_\rho = 0, \Eval 1_\rho = 1, \ldots 
\\[0.5em]
\Eval{\cif{e}{e_1}{e_2}}_\rho = \left \{
\begin{array}{l@{\quad}l}
\Eval{e_1}_\rho, & \mbox{if $\Eval e_\rho = \vtt$} \\
\Eval{e_2}_\rho, & \mbox{if $\Eval e_\rho = \vff$} \\
\bot, & \mbox{otherwise}
\end{array} 
\right . 
\\[1.8em]
\Eval{\cpair{e_1, e_2}}_\rho = (\Eval{e_1}_\rho, \Eval{e_2}_\rho), \quad
\Eval{\cproj_i e}_\rho = v_i \mbox{ (if $\Eval{e}_\rho = (v_1, v_2)$, $i = 1, 2$)}
\\[0.5em]
\Eval{\cinj_1(e)}_\rho = (\mathbb 1, \Eval{e}_\rho) ,
\Eval{\cinj_2(e)}_\rho = (\mathbb 2, \Eval{e}_\rho) 
\\[0.5em]
\Eval{\mathtt{case} \; e \; \mathtt{of} \; \cinj_1 x_1 \; \mathtt{in} \; e_1 ; \cinj_2 x_2 \; \mathtt{in} \; e_2}_\rho = 
\left \{
\begin{array}{l@{\quad}l}
\Eval{e_1}_{\rho[x_1 \mapsto v_1]}, & \mbox{if $\Eval e_\rho = (\mathbb 1, v_1)$} \\
\Eval{e_2}_{\rho[x_2 \mapsto v_2]}, & \mbox{if $\Eval e_\rho = (\mathbb 2, v_2)$} \\
\bot, & \mbox{otherwise}
\end{array} 
\right . 
\\[1.8em]
\Eval{e_1 \ e_2}_\rho = f(v), \mbox{ where $f = \Eval{e_1}_\rho, v = \Eval{e_2}_\rho$}
\\[0.5em]
\Eval{\lambda x : \tau' \ldot e}_\rho = f 
\mbox{ such that \empurple{for every $v \in \Eval{\tau'}$, $f(v) = \Eval{e}_{\rho[x \mapsto v]}$}}
\end{array}
\vspace*{-1.5em}
\]
%% \item An element $v$ of $\Eval \tau$ is \emblue{definable} if there exists a PCF program $e$ of 
%%   type $\tau$ such that $\Eval e = v$.
\item Set model is not sufficient for interpreting recursions: \emred{NOT every function has a fix point}.
\\
PCF model needs more complex structure, e.g.,  \emblue{CPOs}. 
\end{itemize}
\end{slide}

\begin{slide}{Categories}
\begin{itemize}
\item Models of typed $\lambda$-calculi can be viewed with a more abstract and unified structure --- \emred{categories}.
\item A category $\cate$ comprises 
\begin{itemize}\small
\item a collection of \emblue{objects};
\item a collection of \emblue{morphisms} (between objects), and every morphism $f$ has 
  a unique {\em source object} $A$ and {\em target object} $B$, written as $f: A \to B$. 
  We write \emblue{$\cate[A, B]$} for the set of all morphisms from $A$ to $B$;
\item a binary operation \emblue{$\circ$ (composition of morphisms)}: if $f: A \to B$ and  $g: B \to C$ are morphism, 
  then $g \circ f$ is a morphism from $A$ to $C$. \\ 
  $\circ$ is {\em associative}: if $f: A \to B$, $g: B \to C$ and $h: C \to D$ are morphism, 
  then \empurple{$h \circ (g \circ f) = (h \circ g) \circ f$};
\item an \emblue{identity morphism $\pmb{id}_X: X \to X$} (for every object $X$): for every morphism $f: A \to B$, 
  \empurple{$\pmb{id}_B \circ f = f = f \circ \pmb{id}_A$}.
\end{itemize}
\end{itemize}
\end{slide}

\begin{slide}{Category examples}
\begin{itemize}
\item Sets and total functions (\emblue{$\mathcal Set$}) \vspace*{-0.5em}
\begin{itemize}\small
\item sets are obejcts;
\item total functions are morphisms, with their domains and codomains as source and target objects; 
\item $\circ$ is function composition, which is associative;
\item identity function for every set as identity morphisms. 
\vspace*{-1em}
\end{itemize}
\item Sets and partial functions
\item CPOs and continuous functions
\item Groups and group homomorphisms
\item Vector spaces and linear transforms
\item \empurple{Types and terms in typed $\lambda$-calculi}
\end{itemize}
\end{slide}

\begin{slide}{Cartesian closed categories (CCC)}
% More constructions for interpreting typed $\lambda$-calculi:
\begin{itemize}\small
\item \emblue{Terminal object $1$}: for every object $A$, there is 
  \empurple{exactly one morphism (denoted by $!$)} from $A$ to $1$.
$$A \xrightarrow{\qquad ! \qquad} 1$$ \vspace*{-1.5em}
\item \emblue{Product}: a product of two objects $A, B$ is an object \emblue{$A \times B$}, 
  together with two {\em projection morphisms} \emblue{$\pi_1: A \times B \to A, \pi_2 : A \times B \to B$}, 
  such that for every object $C$ and pair of morphisms $f: C \to A, g : C \to B$, 
  there is exactly one {\em mediating morphisms} \emblue{$\cpair{f, g}: C \to A \times B$} 
  making the following diagram commutes, 
  i.e., \empurple{$\pi_1 \circ \cpair{f, g} = f, \pi_2 \circ \cpair{f, g} = g$}:
$$
\xymatrix@+=3pc{
  & C \ar[dl]_{f} \ar[dr]^{g} \ar@{-->}[d]^{\cpair{f, g} }& \\
  A & A \times B \ar[l]_{\quad \pi_1 \quad} \ar[r]^{\quad \pi_2 \quad} & B
}
$$
If $f: A \to C, g: B \to D$ are two morphisms, we write $f \times g : A \times B \to C \times D$ 
for the \emblue{product morphism $\cpair{f \circ \pi_1, g \circ \pi_2}$}.
\end{itemize}
\end{slide}

\begin{slide}{Cartesian closed categories (CCC)}
\begin{itemize}\small
\item \emblue{Exponentiation}: an exponentiation object of two objects $A, B$ is an object $B^A$ 
  together with an {\em evaluation morphism} \emblue{$\pmb{eval}_{AB}: B^A \times A \to B$}
  such that for every object $C$ and morphism $g : C \times A \to B$, 
  there is a unique morphism \emblue{$\pmb{curry}(g) : C \to B^A$} 
  making the following diagram commutes:
$$
\xymatrix@+=3pc{
  C \times A \ar[dr]^{g} \ar@{-->}[d]_{\pmb{curry}(g) \times \pmb{id}_A} & \\
  B^A \times A \ar[r]^{\pmb{eval}_{AB}} & B
}
$$
That is: \empurple{$\pmb{eval}_{AB} \circ (\pmb{curry}(g) \times \pmb{id}_A) = g$}.
\item A \emred{cartesian closed category (CCC)} is a category 
  with a terminal object, products, and exponentiation.
\end{itemize}
\end{slide}

\begin{slide}{Modeling typed $\lambda$-calculli in CCC}
\begin{itemize}
\item Types as objects 
$$\small
\begin{array}{lcl}
\Eval{\tunit} & \defeq & 1 \\
\Eval{\tau_1 \times \tau_2} & \defeq & \Eval{\tau_1} \times \Eval{\tau_2} \\
\Eval{\tau_1 \to \tau_2} & \defeq & \Eval{\tau_2}^{\Eval{\tau_1}}
\end{array}
$$\vspace*{-0.5em}
\item Programs as morphisms
\begin{itemize}\small
\item Typing contexts as products
$$\small
\begin{array}{lcl}
\Eval{\emptyset} & \defeq & 1 \\
\Eval{\Gamma, x : \tau} & \defeq & \Eval\Gamma \times \Eval{\tau}
\end{array}
$$
\end{itemize}
\end{itemize}
\end{slide}

\begin{slide}{Modeling typed $\lambda$-calculli in CCC}
\begin{itemize}
\item Programs as morphisms
\end{itemize}
$$\small
\begin{array}{lcl}
\Eval{\Gamma \vdash * : \tunit} & \defeq & ! : \Eval{\Gamma} \to 1 
\\[0.2em]
\Eval{x_1 : \tau_1, \ldots, x_n : \tau_n \vdash x_i : \tau_i} & \defeq & 
\pi_i : \Eval{\tau_1} \times \cdots \times \Eval{\tau_n} \to \Eval{\tau_i}
\\[0.2em]
\Eval{\Gamma \vdash \lambda x \ldot e : \tau \to \tau'} & \defeq & 
\pmb{curry}(\Eval{\Gamma, x : \tau' \vdash e : \tau}) : \Eval{\Gamma} \to \Eval{\tau}^{\Eval{\tau'}}
\\[0.2em]
\Eval{\Gamma \vdash e_1 \ e_2 : \tau} & \defeq & 
\pmb{eval}_{\Eval{\tau'} \Eval{\tau}} \circ \cpair{\Eval{\Gamma \vdash e_1 : \tau' \to \tau}, \Eval{\Gamma \vdash e_2 : \tau'}}
\\[0.1em]
& & \qquad : \Eval{\Gamma} \to \Eval{\tau} 
\\[0.2em]
\Eval{\Gamma \vdash \cpair{e_1, e_2} : \tau_1 \times \tau_2} & \defeq &
\cpair{\Eval{\Gamma \vdash e_1 : \tau_1}, \Eval{\Gamma \vdash e_2 : \tau_2}} 
: \Eval{\Gamma} \to \Eval{\tau_1} \times \Eval{\tau_2}
\\[0.2em]
\Eval{\Gamma \vdash \cproj_1{e} : \tau_1} & \defeq &
\pi_1 \circ \Eval{\Gamma \vdash e : \tau_1 \times \tau_2} : \Eval{\Gamma} \to \Eval{tau_1}
\\[0.2em]
\Eval{\Gamma \vdash \cproj_2{e} : \tau_2} & \defeq &
\pi_2 \circ \Eval{\Gamma \vdash e : \tau_1 \times \tau_2} : \Eval{\Gamma} \to \Eval{tau_2}
\\[0.2em]
 & \cdots & \cdots
\end{array}
$$
\end{slide}

\begin{slide}{The computational $\lambda$-calculus}
\begin{itemize}
\item A nice framework for \empurple{modeling imperative features in pure functional language} \emgreen{[Moggi'89]}. 
\item Typed $\lambda$-calculus with specfic types for (imperative) computations.
\begin{itemize} \small
\item Types: \vspace*{-0.8em}
$$\tau, \tau', \ldots ::= b \mid \tau \to \tau \mid \ldots \mid \emred{\tcomp \tau} $$
$\tcomp \tau$ is the type for \emblue{computations that can return values of type $\tau$}. 
\item Terms: \vspace*{-0.8em}
$$ 
e, e', \ldots ::= c \mid \lambda x \ldot e \mid e \ e' \mid \ldots \mid \emred{\cval(e)} 
\mid  \emred{\cbind{x}{e}{e'}} 
\vspace*{-0.5em}
$$
$\cval(e)$ is the \emblue{trivial computation} that does nothing but returns $e$ directly; 
$\cbind{x}{e}{e'}$ is the \emblue{sequential computation} that 
executes the computation $e$ first, binds its value to $x$ and executes the computation $e'$.
\item Typing rules:
$$
\prooftree
\Gamma \vdash e : \tau 
\justifies 
\Gamma \vdash \cval(e) : \tcomp \tau 
\endprooftree
\qquad
\prooftree
\Gamma \vdash \emred{e : \tcomp \tau} 
\quad 
\Gamma, \emred{x : \tau} \vdash e' : \tcomp \tau'
\justifies 
\Gamma \vdash \cbind x e {e'} : \tcomp \tau'
\endprooftree
$$
\end{itemize}
\end{itemize}
\end{slide}

\begin{slide}{Modeling imperative computations}
\begin{itemize}\small
\item Non-termination \vspace*{-0.6em} 
\begin{itemize}
\item $\tcomp \tau \defeq \emblue{\tau + \tunit}$: a computation may terminate and return a value of $\tau$, 
  or may not terminate, which is the $\tunit$ case (recall the lifted type that we introduce in PCF store).
\item $\cval(e)$ always teminates and returns the value $e$.
\item $\cbind x e {e'}$ can teminate if both $e$ and $e'$ do, and the final result is the one returned by $e'$; 
  if either of $e, e'$ does not terminate, neither does the entire program. 
\vspace*{-0.8em}
\end{itemize}
\item Exceptions \vspace*{-0.6em}
\begin{itemize}\small
\item $\tcomp \tau \defeq \emblue{\tau + E}$ ($E$ is the type of all exceptions): a computation may return normally 
  or throw an exception. 
\item $\cval(e)$ always returns normal value $e$.
\item $\cbind x e {e'}$: if $e$ returns normally then the program continues with $x$ being replaced by the returned value;
  otherwise, when $e$ throws an exception the program stops with that exception. 
\vspace*{-0.6em}
$$
\prooftree
\Gamma \vdash e: E 
\justifies 
\Gamma \vdash \mathtt{throw}_\tau (e) : \tcomp \tau 
\endprooftree
\qquad
\prooftree
\Gamma \vdash e : \tcomp \tau 
\quad 
\Gamma \vdash e' : E \to \tcomp \tau
\justifies 
\mathtt{catch} \; e \; \mathtt{with} \; e' : \tcomp \tau
\endprooftree
$$
\end{itemize}
\end{itemize}
\end{slide}

\begin{slide}{Modeling imperative computations}
\begin{itemize}\small
\item Stateful computations \vspace*{-0.6em}
\begin{itemize}
\item $\tcomp \tau \defeq \emblue{S \to S \times \tau}$ ($S$ is the type of all states): a computation always 
  executes in a given state and returns a value plus a new state.
\item $\cval(\tau)$ returns the value $e$ and the same state as it receives.
\item $\cbind x e {e'}$ starts the computation $e$ in state $s_0$, which returns a value $v$ and a new state $s_1$, 
  then starts the computation $e''$ in state $s_1$ with $x$ being replaced by $v$.
\item Modeling imperative programs (WHILE) with state monad: \vspace*{-0.5em}
$$
\prooftree
\justifies 
\Gamma \vdash \pmb{valueof} : \tloc \to \emred{\tcomp \tnat}
\endprooftree
\qquad 
\prooftree
\justifies 
\Gamma \vdash \pmb{update} : \tloc \to \tnat \to \emred{\tcomp \tunit}
\endprooftree
$$
Interpretation of expressions are of type $\tcomp \tnat$, which do not change states, 
and commands are $\tcomp \tunit$, which modify states:
$$
\begin{array}{lcl}
\Eval{x := E} & \defeq & \cbind {v} {\Eval E} {\pmb{update}(\Eval{x}, v)} 
\\[0.2em]
\Eval{P_1 ; P_2} & \defeq & \cbind {\_} {\Eval{P_1}} {\Eval{P_2}}
\end{array}
$$
\end{itemize}
\end{itemize}
\end{slide}

\begin{slide}{Functors and natural tranformations}
\begin{itemize}\small
\item A \emblue{functor} $F : \mathcal C \to \mathcal D$ ($\mathcal C, \mathcal D$ are categories) is a map 
taking each $\mathcal C$-object $A$ to a $\mathcal D$-object $F(A)$, and each $\mathcal C$-morphism $f: A \to B$ 
to a $\mathcal D$-morphism $F(f) : F(A) \to F(B)$, such that for all $\mathcal C$-objects $A$ and composable 
$\mathcal C$ morphisms $f$ and $g$, \vspace*{-0.6em}
\begin{itemize}
\item \emblue{$F(\pmb{id}_A) = \pmb{id}_{F(A)}$} and 
\item \emblue{$F(g \circ f) = F(g) \circ F(f)$}. 
\vspace*{-0.6em}
\end{itemize}
Type constructors in programming languages are typical functors. 
\item Let $\mathcal C, \mathcal D$ be categories, and $F, G$ be functors from $\mathcal C$ to $\mathcal D$. 
a \emblue{natural transformation} $\eta$ from $F$ to $G$ (written $\eta: F \to G$) is a function that assigns 
to every $\mathcal C$-object $A$ a $\mathcal D$-morphism $\eta_A: F(A) \to G(A)$, 
such that for every $\mathcal C$-morphism $f: A \to B$, the following diagram commutes in $\mathcal D$:
$$
\xymatrix@+=3pc{
  F(A) \ar[r]^{\eta_A} \ar[d]_{F(f)} &  G(A) \ar[d]^{G(f)} \\
  F(B) \ar[r]^{\eta_B} & G(B)
}
$$
\end{itemize}
\end{slide}

\begin{slide}{Monads}
\begin{itemize}
\item A \emblue{monad} on a category $\mathcal C$ is a tuple $(T, \eta, \mu)$
\begin{itemize}
\item $T: \mathcal C \to \mathcal C$ is a functor from $\mathcal C$ to itself;
\item $\eta: 1_{\mathcal C} \to T$ is a natural transformation from identity functor $1_{\mathcal C}$ to $T$;
\item $\mu: T^2 \to T$ is a natural transformation from functor $T \circ T$ to $T$;
\item the following two diagrams commute:
$$
\xymatrix@+=3pc{
  T^3 \ar[r]^{T \mu} \ar[d]_{\mu_T} &  T^2 \ar[d]^{\mu} \\
  T^2 \ar[r]^{\mu} & T
}
\qquad 
\xymatrix@+=3pc{
  &  T \ar[dl]_{T \eta} \ar@{=}[d] \ar[dr]^{\eta_T} & \\
  T^2 \ar[r]^{\mu} & T & T^2 \ar[l]_{\mu}
}
$$
i.e., $\mu \circ T \mu = \mu \circ \mu_T$ and $\mu \circ T \eta = \mu \circ \eta_T$.
\item A \emblue{strong monad} is also equipped with a morphism 
  \emblue{$\pmb{t}_{A,B}: A \times T B \to T (A \times B)$} for every pair of objects $A, B$.
\end{itemize}
\end{itemize}
\end{slide}

\begin{slide}{Interpreting the computational $\lambda$-calculus}
\begin{itemize}
\item The computational $\lambda$-calculus is interpreted in a CCC with a \emblue{strong monad} 
  $(T, \eta, \mu, \pmb{t})$.
\begin{itemize} \small
\item Computation types
$$\Eval {\tcomp \tau} \defeq T \Eval{\tau} 
\vspace*{-1.5em}
$$
\item Trivial computations \vspace*{-0.6em}
$$
\Eval{\Gamma \vdash \cval(e) : \tcomp \tau} \defeq
\eta_{\Eval \tau} \circ \Eval{\Gamma \vdash e : \tau} 
\vspace*{-1em}
$$ 
$$
\xymatrix@+=3pc{
  \Eval \Gamma \ar[r] ^{\Eval{\Gamma \vdash e: \tau}} & \Eval \tau \ar[r]^{\eta_{\Eval \tau}} & T \Eval \tau 
}
\vspace*{-1.5em}
$$
\item Sequential computations \vspace*{-0.6em}
$$
\begin{array}{l}
\Eval{\Gamma \vdash \cbind x e {e'} : \tcomp \tau'} \defeq 
\\ \qquad\qquad
\mu_{\Eval{\tau'}} \circ T \Eval{\Gamma, x: \tau \vdash e' : \tau'} \circ \pmb{t}_{\Eval\Gamma, \Eval \tau}
\circ \cpair{\pmb{id}_{\Eval \Gamma}, \Eval{\Gamma \vdash e : \tau}}  
\end{array}
$$
$$\small
\xymatrix@C=2.5pc@R=1.5pc{
  \Eval \Gamma \ar[d]_{\cpair{\pmb{id}, \Eval{\Gamma \vdash e: \tau}}} & & & T \Eval{\tau'} 
  \\
  \Eval \Gamma  \times T \Eval \tau \ar[r]^(.48){\pmb{t}_{\Eval \Gamma, \Eval \tau}} 
  & T (\Eval \Gamma \times \Eval \tau) \ar[rr]^(.52){T \Eval{\Gamma, x: \tau \vdash e': \tcomp \tau'}} 
  & & T T \Eval{\tau'} \ar[u]_{\mu_{\Eval{\tau'}}} 
}
$$
\end{itemize}
\end{itemize}
\end{slide}

\end{document}
 

